\documentclass{mdtfurb}
\usepackage[utf8]{inputenc}
\usepackage[T1]{fontenc}
\usepackage{geometry}
\usepackage{graphicx}
\usepackage{float}
\usepackage{hyperref}
\usepackage[english,indonesian]{babel}

\title{Optimasi Pemrosesan Template LaTeX Menggunakan Pendekatan Berbasis AI}
\author{Academic Support Agent}
\date{\today}

\begin{document}
\maketitle

\begin{abstract}

\end{abstract}

\chapter{PENDAHULUAN}


\chapter{TINJAUAN PUSTAKA}


% Original content (commented out):
% %=========================================================================
%                       AVISOS - LEIA COM ATENÇÃO!
% 1) A Biblioteca Universitária não fornece suporte à linguagem LaTeX.
% 2) Este modelo pode ser baixado no formato .zip para ser editado em outros editores LaTeX para desktop. Para isso, basta acessar o Menu no canto superior direito e clicar em "Download Source" para baixar o arquivo .zip
%=========================================================================


%%----------------INÍCIO--PREÂMBULO--------------------------------
%%ATENÇÃO!!!! TENHA MUITO CUIDADO AO MODIFICAR O PREÂMBULO, SÓ O MODIFIQUE SE TIVER CERTEZA DO QUE ESTÁ FAZENDO!!!!

%\listfiles
\documentclass[diss]{mdtfurb}
% um tipo específico de monografia pode ser informado como parâmetro opcional: tese, diss, tcc, espec
%EXEMPLO: \documentclass[tese]{mdtfurb}

% para gerar uma versão frente-e-verso, use a opção 'twoside' e 'openright':
%EXEMPLO: \documentclass[diss,twoside,openright]{mdtfurb}

\usepackage[T1]{fontenc}        % pacote para conj. de caracteres correto
\usepackage{fix-cm} %para funcionar corretamente o tamanho das fontes da capa
\usepackage{times, color, xcolor}       % pacote para usar fonte Adobe Times e cores
\usepackage[utf8]{inputenc}   % pacote para acentuação
\usepackage{graphicx}  % pacote para importar figuras
\usepackage{amsmath,latexsym,amssymb} %Pacotes matemáticos
\usepackage[hidelinks, 
            bookmarksopen=true,linktoc=none,colorlinks=true,
            linkcolor=black,citecolor=black,filecolor=magenta,urlcolor=black,
            pdftitle={Título da Dissertação ou Trabalho ....},
            pdfauthor={Nome Autor Sobrenome},
            pdfsubject={Dissertação de Mestrado},
            pdfkeywords={Dissertação, Modelo, LaTeX}
            ]{hyperref} %hidelinks disponível no pacote hyperref a partir da versão 2011-02-05  6.82a
%Nesse caso, hidelinks retira os retângulos em volta dos links das referências
    \urlstyle{same} %Links aparecem com a mesma fonte do texto
    
%Margens:
\usepackage[inner=30mm,outer=20mm,top=30mm,bottom=20mm]{geometry} 

%Configuração das referências:
%Atenção: utilizar Biber como compilador de referências para funcionar corretamente
\usepackage[style=abnt, noslsn]{biblatex} %configura referências --> remover "bftitles" caso esteja usando TeXLive 2017 ou superior (verificar no arquivo de Log)
    \addbibresource{referencias/referencias.bib} %localização do arquivo de referências
    
    
\usepackage{csquotes} %configuração de idioma   

\usepackage{enumitem} %configuração de listas

\usepackage{chngcntr} %numeração contínua de figuras, tabelas e equações
	\counterwithout{figure}{chapter}
	\counterwithout{table}{chapter}
	\counterwithout{cquadro}{chapter}
	\counterwithout{equation}{chapter}


%==============================================================================
% Se o pacote hyperref foi carregado a linha abaixo corrige um bug na hora
% de montar o sumário da lista de figuras e tabelas
%==============================================================================
\input{macros/bugcaption}

%%------------FINAL--PREÂMBULO------------------------------------


%==============================================================================
% Identificação do trabalho
%==============================================================================
\title{Título do trabalho: Subtítulo}

\author{Sobrenome}{Nome Autor}
%Descomentar se for uma "autora"
%\autoratrue

\course{Bacharelado/Licenciatura/Programa de Pós-Graduação em...} %Nome do curso
\altcourse{Curso de... ou Programa de Pós-Graduação em...} %Nome do curso (folha de rosto)

\institute{Centro de...} %Nome do centro
\degree{Bacharel/Licenciado/Tecnólogo/Mestre/Doutor em...} %Grau obtido

%Orientador
\advisor[Prof.]{Dr.}{Silva}{João da}
%Se for uma ``orientadora'' descomentar a linha baixo
%\orientadoratrue

%Co orientador, comentar se não existir
\coadvisor[Prof.]{Drª.}{Santos}{Maria dos}
\coorientadoratrue %Se for uma ``Co-Orientadora''


%Palavras chave
\keyword{Dissertação} 
\keyword{Modelo}
\keyword{LaTeX}

%%=============================================================================
%% Início do documento
%%=============================================================================
\linespread{1.0}
\begin{document}

%%=============================================================================
%% Capa e folha de rosto
%%=============================================================================
\maketitle

%%=============================================================================
%% Catalogação  e Folha de aprovação
%%=============================================================================

%Ficha catalográfica gerada pela biblioteca, somente para mestrado e doutorado. A ficha deve ser solicitada pelo e-mail protec@furb.br quando o trabalho estiver totalmente concluído. Comentar caso não haja
\includeCIP{CIP.pdf}

%Folha de aprovação: escanear a versão assinada e incluir dentro do projeto no formato PDF. Comentar caso não haja
\includeaprov{aprov.pdf}


%%=============================================================================
%% Dedicatória (opcional)
%%=============================================================================
\chapter*{}

\begin{flushright}
\mbox{}\vfill
{\sffamily\itshape À FURB ......}
\end{flushright}

%%=============================================================================
%% Agradecimentos
%%=============================================================================

\chapter*{Agradecimentos} 
Obrigado ao \LaTeX por facilitar a digitação do trabalho


%%=============================================================================
%% Epígrafe (opcional)
%%=============================================================================
\chapter*{}

\begin{flushright}
\mbox{}\vfill
{\sffamily\itshape
``Frase da epígrafe'' \\ }
--- \textsc{SOBRENOME, 2019}
\end{flushright}


%%=============================================================================
%% Resumo
%%=============================================================================
\begin{abstract}
Este é o resumo do trabalho ... Este é o resumo do trabalho ...
Este é o resumo do trabalho ... Este é o resumo do trabalho ...
Este é o resumo do trabalho ... Este é o resumo do trabalho ...
Este é o resumo do trabalho ... Este é o resumo do trabalho ...
Este é o resumo do trabalho ... Este é o resumo do trabalho ...
Este é o resumo do trabalho ... Este é o resumo do trabalho ...
\end{abstract}

%%=============================================================================
%% Abstract
%%=============================================================================
% resumo na outra língua
% como parametros devem ser passados o titulo, o nome do curso,
% as palavras-chave na outra língua, separadas por pontos
%o mês em inglês
%o a sigla do dia em inglês: st, nd, th ...
\begin{englishabstract}
{Dissertation Title}
{Post-Graduate Program in Chemistry}
{Keywords1. Keyword2}
{March}
{st}
Abstract ... Abstract ...Abstract ...Abstract ...Abstract ...Abstract ...Abstract ...
Abstract ...Abstract ...Abstract ...Abstract ...Abstract ...Abstract ...Abstract ...
Abstract ...Abstract ...Abstract ...Abstract ...Abstract ...Abstract ...Abstract ...
Abstract ...Abstract ...Abstract ...Abstract ...Abstract ...Abstract ...Abstract ...
\end{englishabstract}

%% Lista de Ilustrações (opc)
%% Lista de Símbolos (opc)
%% Lista de Anexos e Apêndices (opc)

%%=============================================================================
%% Lista de ilustrações (comentar se não houver)
%%=============================================================================
\listoffigures

%%=============================================================================
%% Lista de tabelas (comentar se não houver)
%%=============================================================================
\listoftables

%%=============================================================================
%% Lista de Apêndices (comentar se não houver)
%%=============================================================================
%\listofappendix

%%=============================================================================
%% Lista de Anexos (comentar se não houver)
%%=============================================================================
%\listofannex

%%=============================================================================
%% Lista de abreviaturas e siglas
%%=============================================================================
 %o parametro deve ser a abreviatura mais longa
\begin{listofabbrv}{UbiComp}
   \item [BNF] \textit{Backus-Naur Form}
   \item [UbiComp] Computação Ubíqua
\end{listofabbrv}


%%=============================================================================
%% Lista de simbolos (opcional)
%%=============================================================================
%Simbolos devem aparecer conforme a ordem em que aparecem no texto
% o parametro deve ser o símbolo mais longo
\begin{listofsymbols}{teste}
  \item [$\varnothing$] vazio
  \item [$\Gamma$]  Gama
  \item [$\forall$] Para todo
\end{listofsymbols}

%%=============================================================================
%% Sumário
%%=============================================================================
\tableofcontents


%%=============================================================================
%% Início da dissertação
%%=============================================================================
\setlength{\baselineskip}{1.5\baselineskip}

%Adiciona cada capitulo
\include{capitulos/introducao}
\include{capitulos/desenvolvimento}
\include{capitulos/consideracoes_finais}

\setlength{\baselineskip}{\baselineskip}

%%=============================================================================
%% Referências
%%=============================================================================

\begingroup
\setlength\bibitemsep{12pt}
\printbibliography[heading=bibintoc,title={Refer{\^e}ncias}]
\endgroup


%%=============================================================================
%% Apêndices
%%=============================================================================

\include{capitulos/apendicea}
\include{capitulos/apendiceb}




%%=============================================================================
%% Anexos
%%=============================================================================

\include{capitulos/anexoa}

\end{document}


\end{document}